\section{Experimental details}
\label{app:experimental-details}

\subsection{Datasets and evaluation protocols}

\placeholder{Report dataset splits, preprocessing, reasoning depth, selectors,
the distinction between $K_{\mathrm{train}}$ and $K_{\mathrm{eval}}$, evaluation
seeds, and confidence-interval construction for Sudoku-Extreme, Maze-Hard,
and ARC-AGI-2.}

\paragraph{Maze-Hard candidate-coverage protocol.}
We evaluate FB, ordinary IVON, PTRM, and W-PTRM on the same 1,000-problem test
split with the TRM backbone, reasoning depth 16, and candidate budgets up to
ten. FB, PTRM, and W-PTRM use ordered prefixes from saved $K_{\mathrm{eval}}=10$
banks under seed-0 reset sampling. The reported IVON values use the updated
evaluation at $K\in\{1,2,4,6,8,10\}$. The earlier IVON bank produced on
August 26, 2026 is deprecated and excluded from the reported comparison.

\subsection{Baseline details}

\paragraph{Deterministic reasoning models.}
\textbf{HRM} couples recurrent modules that operate at different time
scales~\citep{wang2025hierarchical}. \textbf{TRM} replaces the hierarchy with
one small recursive network~\citep{jolicoeurmartineau2025trm}.
\textbf{FPRM} uses fixed-point convergence to halt a looped
Transformer~\citep{movahedi2026fixedpoint}.

\paragraph{Sampling-based reasoning methods.}
\textbf{GRAM} models reasoning as stochastic latent trajectories and samples
several trajectories in parallel~\citep{baek2026generative}. \textbf{PTRM}
keeps the model parameters fixed, injects Gaussian noise into the latent state
at each deep-recursion step, and selects among the resulting
trajectories~\citep{sghaier2026probabilistic}. \textbf{W-PTRM} and
\textbf{W-FPRM} are non-learned baselines that sample each candidate's
parameters from a fixed Gaussian centered at the released TRM or FPRM
checkpoint. They add no latent noise.

\paragraph{Posterior optimization methods.}
\textbf{IVON} (Improved Variational Online Newton) maintains a sampleable
diagonal Gaussian through an online variational update and optimizes the
single-draw posterior objective~\citep{shen2024variational}. \textbf{FB
(Ours)} uses the same IVON posterior parameterization and initialization, but
optimizes the finite-$K$ objective with Fenchel--Bregman updates. For learned
posterior methods, we initialize the posterior mean from a released task
checkpoint and fine-tune its mean and diagonal precision on the task training
split. Fine-tuning adapts the solver distribution to the chosen objective,
while W-PTRM and W-FPRM isolate sampling without this adaptation.

\subsection{Training and compute accounting}

\placeholder{Report optimization hyperparameters, posterior and dual-state
initialization, training-update budgets, hardware, batch sizes, mixed-precision
settings, solver-call accounting, wall-clock measurement, and peak-memory
measurement.}

\paragraph{Ordinary-IVON baseline.}
We report the updated IVON Pass@$K$ evaluation designated on August 27, 2026.
The previous local HPO, checkpoint, and candidate bank are deprecated and are
not used for the reported curve.

